\section*{Preface}

{\fontsize{10.5}{15}\selectfont
This book, titled \emph{Preparation Course for PyTorch} (PCPT), grew directly out of our teaching experience over the past several years. It reflects our experience working with students in lab courses, research internships, and thesis projects at the International Audio Laboratories Erlangen, a joint institution of the Friedrich-Alexander-Universität Erlangen-Nürnberg~(FAU) and the Fraunhofer Institute for Integrated Circuits IIS. Much of this work is situated in audio, music, and speech processing, with a focus on one-dimensional signals and their role in bridging classical signal processing with deep learning and data-driven methods in an interdisciplinary context.

In this setting, we repeatedly observed a similar pattern. Many of our Master's students in engineering, computer science, and related fields come with complementary backgrounds: some have a solid foundation in signal processing but lack confidence when using modern programming frameworks such as PyTorch, while others are more experienced in programming but less familiar with the underlying signal processing concepts. As a result, the connection between theory and practice is often incomplete. Students with a signal processing background are comfortable with concepts such as convolution, filtering, or recursion, but may struggle to translate these ideas into code. Conversely, students with stronger programming experience can apply practical tools effectively, yet sometimes lack a deeper understanding of the underlying mathematical principles. The link between theory and practice therefore remains fragile in both directions. At the same time, standard programming courses are often too generic for our needs, while many machine learning resources focus on image data and do not align well with a signal processing perspective. In our study programs, programming skills are frequently assumed at the Master's level without being systematically taught, and we often found ourselves addressing these gaps within lab courses, internships, and student projects.

To address this situation more systematically, we developed the \emph{Preparation Course for Python} (PCP)\footnote{\url{https://audiolabs-erlangen.de/PCP}}, which introduces programming concepts in Python with a focus on basic digital signal processing~(DSP) topics such as signals, sampling, exponential functions, and the discrete Fourier transform. Building on this foundation, PCPT extends the approach towards PyTorch and modern machine learning. The guiding idea is to start from familiar concepts such as one-dimensional signals and time-series data, and to gradually introduce programming and learning-based methods. Over time, the PCPT course evolved into something we had not originally planned. It is neither a pure programming course nor a conventional machine learning introduction, but an attempt to build a conceptual bridge between signal processing and data-driven modeling.

From the beginning, we relied on Jupyter notebooks as the main teaching medium.\footnote{\url{https://audiolabs-erlangen.de/PCPT}} Their interactive nature allows us to combine explanations, equations, visualizations, and executable code in a single environment, encouraging experimentation and exploration. Many parts of the material emerged directly from such hands-on sessions, often shaped by questions and discussions with students. At the same time, we found that working exclusively with notebooks can make it difficult to follow a clear and structured narrative. Introducing an additional medium such as a printed text provides a complementary perspective. While notebooks offer a local, interactive, and execution-driven environment, a book provides a more global, linear, and reflective view of the material. This combination of digital and analog formats supports different modes of learning, allowing readers to step back, annotate, and revisit the content more easily. This insight led to the creation of this book. We do not see it as a replacement for the notebooks, but as a complementary view of the same material. The book offers a coherent and structured presentation, while the notebooks remain the place for experimentation. In practice, we recommend alternating between both: reading, implementing, and reflecting. This interplay has proven particularly effective in our teaching.

Technically, the book is a direct transformation of the notebooks. Each notebook is cleaned and executed in a reproducible way, then converted into a typeset format where text, figures, and code are consistently integrated. The individual units are merged into a single document with unified layout and typography. While the visual appearance may differ slightly from the notebooks, the content remains identical.

In preparing this material, we deliberately focus on fundamental ideas rather than rapidly changing technical details. The examples are intentionally simple, typically based on synthetic data, using reduced sampling rates and only a small number of training and test examples, while still being motivated by real-world scenarios in audio and music processing. This allows us to present concepts in a controlled and transparent way without losing the connection to practical applications. All examples are designed to run on standard hardware, so no specialized GPU is required; a standard notebook computer is sufficient. This enables quick experimentation and keeps the focus on conceptual understanding rather than computational overhead. Many parts of this book have been shaped by interactions with students. Their questions have helped us refine explanations and rethink how concepts should be introduced. In this sense, the book reflects not only our perspective as instructors, but also the learning process of the students.

This book is published through FAU~University Press, with a digital version freely available in open access. At the same time, it can be ordered worldwide as a professionally printed volume. We deliberately chose this dual format to combine broad accessibility with the benefits of a physical book that can be used alongside the interactive notebooks.

This work was funded by the Deutsche Forschungsgemeinschaft (DFG, German Research Foundation) under Grant No.~500643750 (MU~2686/15-1). The material presented in this book builds on insights and contributions from many collaborators and students. We would like to thank former and current members of our group, as well as collaborators and colleagues, for their support and feedback, in particular (in alphabetical order): Stefan Balke, Michael Krause, Yigitcan Özer, Sebastian Rosenzweig, Simon Schwär, Stefan Turowski, Frank Zalkow, and Johannes Zeitler.

We hope that this book, together with the accompanying notebooks, provides a useful pathway to learning PyTorch and modern machine learning, grounded in signal processing and supported by hands-on experimentation.

\vspace{2em}

\noindent
Erlangen, April 2026\\
Meinard Müller and Sebastian Strahl
}